\documentclass[11pt]{article}

% ACL 2026 final submission style
\usepackage[final]{acl}

% Standard package includes
\usepackage{times}
\usepackage{latexsym}
\usepackage[T1]{fontenc}
\usepackage[utf8]{inputenc}
\usepackage{microtype}
% \usepackage{inconsolata}

% Math and Symbols
\usepackage{amsmath,amssymb,amsfonts,mathtools,bm}

% Graphics, TikZ and Colors
\usepackage{graphicx}
\usepackage{xcolor}
\usepackage{tikz}
\usetikzlibrary{positioning,shapes.geometric,arrows.meta,calc,decorations.pathreplacing}

% Tables and Formatting
\usepackage{booktabs}
\usepackage{multirow}
\usepackage{tabularx}
\usepackage{array}
\usepackage{enumitem}
\usepackage{listings}

% Color Definitions
\definecolor{navyblue}{RGB}{0,32,96}
\definecolor{darkgreen}{RGB}{0,100,0}
\definecolor{darkred}{RGB}{140,20,20}
\definecolor{codegray}{rgb}{0.5,0.5,0.5}
\definecolor{backcolour}{rgb}{0.96,0.96,0.96}

% Listings styling
\lstdefinestyle{promptStyle}{
    backgroundcolor=\color{backcolour},
    commentstyle=\color{darkgreen},
    keywordstyle=\color{navyblue}\bfseries,
    numberstyle=\tiny\color{codegray},
    stringstyle=\color{darkred},
    basicstyle=\ttfamily\scriptsize,
    breaklines=true,
    captionpos=b,
    keepspaces=true,
    showstringspaces=false,
    tabsize=2,
    aboveskip=4pt,
    belowskip=4pt
}
\lstset{style=promptStyle}

\setlength\titlebox{4.3cm}

\title{Autonomous Knowledge Graph Exploration via ARK: Fine-Tuning and\\Distilling Small Language Models for Adaptive Retrieval}

\author{
  \begin{tabular}{ccc}
    \textbf{Nguyen Hoai Nam} & \textbf{Pham Cung Le Nhan Vu} & \textbf{Nguyen Khoa Tan Tien} \\
    Student ID: 25C01014 & Student ID: 25C01049 & Student ID: 25C01048
  \end{tabular} \\[0.35cm]
  VNUHCM -- University of Science \\
  Faculty of Information Technology \\
}

\begin{document}
\maketitle

% =============================================================================
% ABSTRACT
% =============================================================================
\begin{abstract}
Retrieval-Augmented Generation (RAG) over heterogeneous, text-rich Knowledge Graphs (KGs) demands a dynamic equilibrium between global discovery across unconnected entities and multi-hop relational traversal. Conventional vector-based similarity methods suffer from relational blindness and shallow coverage, whereas rigid multi-hop algorithms rely on fragile seed selection and fixed hop limits. In this paper, we study the \textbf{ARK} (\textbf{A}daptive \textbf{R}etriever of \textbf{K}nowledge) framework---an autonomous multi-agent exploration system operating over a minimal two-tool action space (\textit{Global Search} via BM25 and \textit{Neighborhood Exploration} via one-hop expansion) aggregated through \textbf{Voting Rank Fusion} (VRF). To eliminate dependency on proprietary frontier models during multi-step graph exploration, we design a comprehensive trajectory mining and knowledge distillation pipeline. We construct Oracle shortest-path trajectories for Supervised Fine-Tuning (SFT) and harvest exploratory trajectories from Gemini~3.6~Flash filtered by Rejection Sampling ($\text{Recall@20} > 0$). We fine-tune \texttt{Qwen3-4B} via full-precision LoRA and distill \texttt{Qwen3-0.6B} via QLoRA (INT4 NF4). On the biomedical \textbf{STaRK-PrimeKG} benchmark (129k nodes, 8.1M edges), fine-tuned \texttt{Qwen3-4B} achieves $\text{Hit@1} = 0.2000$ with a single agent, matching the 5-agent ensemble of the teacher model ($\text{Hit@1} = 0.2060$), while student models converge rapidly within 111--120 optimization steps. Our findings demonstrate that lightweight, fine-tuned language models can master complex stochastic graph exploration while dramatically reducing inference costs.
\end{abstract}

% =============================================================================
% 1. INTRODUCTION
% =============================================================================
\section{Introduction}
\label{sec:intro}

Retrieval-Augmented Generation (RAG) has emerged as the standard paradigm for grounding Large Language Models (LLMs) in verified external knowledge \citep{lewis2020rag,pan2024unifying}. While unstructured text retrieval has seen substantial advances, structured and heterogeneous Knowledge Graphs (KGs) offer explicit, verifiable relational facts that are critical for domains such as biomedicine, scientific discovery, and multi-hop question answering \citep{chandak2023primekg,wu2024stark}.

\begin{figure*}[t]
\centering
\resizebox{0.98\textwidth}{!}{%
\begin{tikzpicture}[
    node distance=0.6cm and 0.9cm,
    querybox/.style={rectangle, draw=black, fill=gray!12, thick, rounded corners=3pt, align=center, inner sep=6pt, font=\small},
    agentbox/.style={rectangle, draw=darkred!90!black, fill=red!6, thick, rounded corners=3pt, align=center, inner sep=5pt, font=\small},
    toolbox/.style={rectangle, draw=navyblue, fill=blue!6, thick, rounded corners=3pt, align=center, inner sep=5pt, font=\small},
    fusionbox/.style={rectangle, draw=darkgreen!90!black, fill=green!6, thick, rounded corners=3pt, align=center, inner sep=6pt, font=\small},
    outbox/.style={rectangle, draw=orange!90!black, fill=yellow!18, thick, rounded corners=3pt, align=center, inner sep=6pt, font=\small},
    arrow/.style={-Latex, thick, navyblue},
    toolarrow/.style={-Latex, dashed, thick, darkred!80!black}
]
% Query Node
\node[querybox] (query) {\textbf{User Query} $q$ \\ \textit{``What protein interacts with Target $X$}\\\textit{and is indicated for Disease $Y$?''}};

% Parallel Agents
\node[agentbox, above right=-0.3cm and 0.9cm of query] (agent1) {\textbf{Agent 1} ($\tau_1$)\\\textit{Stochastic Trajectory ($T=0.7$)}};
\node[agentbox, below=0.35cm of agent1] (agent2) {\textbf{Agent 2} ($\tau_2$)\\\textit{Stochastic Trajectory ($T=0.7$)}};
\node[agentbox, below=0.35cm of agent2] (agent3) {\textbf{Agent $M$} ($\tau_M$)\\\textit{Stochastic Trajectory ($T=0.7$)}};

% Knowledge Graph & Tools
\node[toolbox, right=1.0cm of agent2] (kg) {\textbf{Heterogeneous Knowledge Graph $\mathcal{G}$}\\\textbf{Minimal 2-Tool Action Space $\mathcal{A}$}\\[3pt]
\begin{tabular}{l}
$\bullet$ \textbf{GlobalSearch}($q_{\text{sub}}, k$): BM25 Lexical Node Discovery \\
$\bullet$ \textbf{NeighborhoodExploration}($v, q_{\text{sub}}, k$): 1-Hop Traversal \\
$\bullet$ \textbf{AddToAnswer}($v_1, \dots, v_m$): Accumulate Candidates $\mathcal{S}_m$ \\
$\bullet$ \textbf{Finish}(): Terminate Trajectory
\end{tabular}};

% Voting Rank Fusion
\node[fusionbox, right=1.0cm of kg] (vrf) {\textbf{Voting Rank Fusion (VRF)}\\[2pt]
$\text{Score}_{\text{VRF}}(v) = \sum\limits_{m=1}^M \frac{\mathbb{I}(v \in \mathcal{S}_m)}{\text{rank}_m(v) + \kappa}$\\[4pt]
\textit{Aggregates candidate sets $\{\mathcal{S}_1, \dots, \mathcal{S}_M\}$}};

% Final Output
\node[outbox, right=0.8cm of vrf] (output) {\textbf{Ranked Candidate Set}\\[2pt]\textbf{Top-$K$ Entities} $\mathcal{R}_K$\\[3pt]\textit{Ground-Truth Targets $\mathcal{V}^*$}};

% Connections
\draw[arrow] (query.east) -- +(0.4,0) |- (agent1.west);
\draw[arrow] (query.east) -- (agent2.west);
\draw[arrow] (query.east) -- +(0.4,0) |- (agent3.west);

\draw[toolarrow, <->] (agent1.east) -- (kg.north west);
\draw[toolarrow, <->] (agent2.east) -- (kg.west);
\draw[toolarrow, <->] (agent3.east) -- (kg.south west);

\draw[arrow] (kg.east) -- node[above, font=\scriptsize, align=center]{Candidate Sets\\$\{\mathcal{S}_1, \dots, \mathcal{S}_M\}$} (vrf.west);
\draw[arrow] (vrf.east) -- (output.west);

\end{tikzpicture}%
}
\caption{The ARK Multi-Agent Autonomous Exploration Architecture. Multiple independent stochastic agents execute dual-action trajectories ($\text{GlobalSearch}$ for breadth discovery and $\text{NeighborhoodExploration}$ for relational depth) over the heterogeneous graph $\mathcal{G}$. Proposed candidate sets $\mathcal{S}_m$ are aggregated into the final ranking $\mathcal{R}_K$ via Voting Rank Fusion.}
\label{fig:ark_arch}
\end{figure*}

However, retrieving relevant subgraphs or target entities from large-scale, text-rich KGs presents two fundamental challenges:
\begin{enumerate}[leftmargin=*]
    \item \textbf{Shallow Coverage vs. Relational Blindness:} Dense semantic vector retrieval indexes textual node descriptors independently. While effective for initial lexical matching, it cannot follow relational paths ($A \xrightarrow{r_1} B \xrightarrow{r_2} C$), failing multi-hop queries where the answer node shares low direct semantic similarity with the question.
    \item \textbf{Brittleness of Fixed Graph Traversal:} Graph traversal methods typically depend on exact seed entity linking. If the initial seed identification fails, or if the reasoning path spans across disconnected graph components, traversal halts prematurely.
\end{enumerate}

To overcome these dual bottlenecks, \citet{polonuer2026ark} introduced \textbf{ARK} (\textbf{A}daptive \textbf{R}etriever of \textbf{K}nowledge), an agentic retrieval framework that equips an LLM with a minimal two-tool interface: \textit{Global Search} (lexical BM25 search over node text descriptors) for breadth-oriented discovery, and \textit{Neighborhood Exploration} (one-hop neighbor expansion ranked by BM25) for depth-oriented relational traversal. By dynamically interleaving these operations, ARK autonomously navigates the graph without fixed hop bounds or rigid seed dependencies.

Despite its strong performance, running multi-agent stochastic rollouts using proprietary commercial LLMs (e.g., GPT-4 or Gemini 3.6 Flash) incurs high latency and substantial API costs, making production deployment prohibitive.

In this work, we address the challenge of \textbf{small model specialization for agentic graph exploration}. We develop an end-to-end trajectory mining and parameter-efficient fine-tuning/distillation pipeline that transfers graph exploration capabilities from Oracle shortest-paths and frontier teacher models to lightweight, open-source language models (\texttt{Qwen3-4B} and \texttt{Qwen3-0.6B}).

Our core contributions are as follows:
\begin{itemize}[leftmargin=*]
    \item \textbf{Formalization of Autonomous Graph Navigation:} We formalize the ARK multi-agent state space and its Voting Rank Fusion (VRF) aggregation mechanism across stochastic parallel trajectories.
    \item \textbf{Dual Trajectory Mining Pipeline:} We design an Oracle shortest-path trajectory generator and an exploratory teacher distillation pipeline with \textit{Rejection Sampling} ($\text{Recall@20} > 0$), filtering noisy exploration traces into high-quality training demonstrations.
    \item \textbf{Parameter-Efficient Specialization:} We formulate an assistant-masked Supervised Fine-Tuning (SFT) objective and implement full-precision LoRA for \texttt{Qwen3-4B} and INT4 NF4 QLoRA for \texttt{Qwen3-0.6B}.
    \item \textbf{Comprehensive Empirical Analysis:} On the challenging biomedical \textbf{STaRK-PrimeKG} benchmark, we demonstrate that a single fine-tuned \texttt{Qwen3-4B} agent matches the retrieval performance of a 5-agent Gemini 3.6 Flash teacher ($\text{Hit@1} = 0.2000$ vs. $0.2060$). We further present exhaustive tool ablations, convergence analyses, and trajectory scaling behaviors.
\end{itemize}

% =============================================================================
% 2. RELATED WORK
% =============================================================================
\section{Related Work}
\label{sec:related}

\paragraph{Knowledge Graph Question Answering and Graph RAG.}
Traditional KGQA methods translate natural language into structured SPARQL or Cypher queries \citep{sun2023thinkongraph} or rely on Graph Neural Networks (GNNs) for node classification \citep{yasunaga2021qa}. Recent Graph RAG paradigms construct global text-derived knowledge graphs and summarize communities hierarchically \citep{edge2024graphrag}. However, on text-rich heterogeneous KGs where node attributes contain rich clinical or domain documentation, static query generation or flat vector chunking fails to capture multi-hop semantics \citep{wu2024stark}.

\paragraph{Agentic Tool Use for Information Retrieval.}
Equipping language models with external search and calculator tools has proven effective for complex reasoning \citep{yao2022react,chen2023teaching}. While general-purpose agents execute open-ended web browsing, graph exploration demands domain-specific inductive biases: balancing global keyword lookup with relational neighborhood expansion. ARK \citep{polonuer2026ark} formalizes this trade-off into a two-operation action space.

\paragraph{Knowledge Distillation and Trajectory Fine-Tuning.}
Distilling reasoning capabilities from large frontier teachers into compact models via imitation learning has shown strong promise \citep{ouyang2022instructgpt}. In agentic domains, Rejection Sampling Fine-Tuning (RFT) filters successful execution traces to prevent student models from internalizing hallucinated or sub-optimal tool calls. We extend trajectory distillation to multi-hop graph retrieval using parameter-efficient adapters \citep{hu2022lora,dettmers2023qlora}.

% =============================================================================
% 3. THE ARK EXPLORATION FRAMEWORK
% =============================================================================
\section{The ARK Exploration Framework}
\label{sec:framework}

\subsection{Problem and Graph Formulation}
Let $\mathcal{G} = (\mathcal{V}, \mathcal{E}, \mathcal{T}_V, \mathcal{T}_E)$ denote a heterogeneous, text-rich Knowledge Graph, where $\mathcal{V}$ is the entity node set, $\mathcal{E} \subseteq \mathcal{V} \times \mathcal{T}_E \times \mathcal{V}$ is the typed directed relationship set, and $\mathcal{T}_V, \mathcal{T}_E$ represent node and edge type alphabets. Each node $v \in \mathcal{V}$ is associated with a textual descriptor $\mathcal{D}(v)$, which concatenates its canonical name, type label, and descriptive attributes.

Given a user query $q$, the objective is to retrieve a ranked list of candidate answer nodes $\mathcal{R}_K = [v_{(1)}, v_{(2)}, \dots, v_{(K)}] \subset \mathcal{V}$ that contain the ground-truth target entities $\mathcal{V}^* \subset \mathcal{V}$.

\subsection{Minimal Action Space $\mathcal{A}$}
ARK equips the agent with an action space $\mathcal{A}$ consisting of four primitive operations: $\text{GlobalSearch}$, $\text{NeighborhoodExploration}$, $\text{AddToAnswer}$, and $\text{Finish}$:

\begin{enumerate}[leftmargin=*]
    \item $\text{\textbf{GlobalSearch}}(q_{\text{sub}}, k)$: Performs lexical BM25 matching \citep{robertson2009bm25} over all descriptors $\{\mathcal{D}(v) \mid v \in \mathcal{V}\}$, returning the top-$k$ most relevant node IDs:
    \begin{equation}
    \label{eq:bm25}
    \resizebox{0.90\columnwidth}{!}{$
    \displaystyle \text{BM25}(q_{\text{sub}}, v) = \sum_{t \in q_{\text{sub}}} \text{IDF}(t) \cdot \frac{\text{TF}(t, \mathcal{D}(v)) \cdot (k_1 + 1)}{\text{TF}(t, \mathcal{D}(v)) + k_1 \cdot \left(1 - b + b \cdot \frac{|\mathcal{D}(v)|}{\text{avgdl}}\right)}
    $}
    \end{equation}
    where $\text{TF}(t, \mathcal{D}(v))$ denotes the term frequency of token $t$ in descriptor $\mathcal{D}(v)$, and $\text{avgdl}$ is the average descriptor length across the graph. This enables discontinuous jumps across the graph topology to locate seed entities matching sub-clauses of $q$.

    \item $\text{\textbf{NeighborhoodExploration}}(v, q_{\text{sub}}, k)$: Expands the 1-hop relational neighborhood $\mathcal{N}(v) = \{u \in \mathcal{V} \mid (v, r, u) \in \mathcal{E} \lor (u, r, v) \in \mathcal{E}\}$. Adjacent candidate nodes are ranked according to the BM25 score of their descriptors against $q_{\text{sub}}$, returning the top-$k$ neighbors along with their connecting relation types $r$. This provides depth-oriented relational traversal.

    \item $\text{\textbf{AddToAnswer}}(v_1, \dots, v_m)$: Appends identified candidate nodes $v_i$ into the agent's internal ordered answer buffer $\mathcal{S} = [v_1, v_2, \dots]$.

    \item $\text{\textbf{Finish}}()$: Terminates the trajectory when the agent determines sufficient evidence has been gathered.
\end{enumerate}

Algorithm~\ref{alg:ark_rollout} outlines the single-agent decision loop across exploration steps.

\begin{figure}[t]
\centering
\begin{lstlisting}[mathescape=true, caption={ARK Single-Agent Trajectory Rollout}, label={alg:ark_rollout}]
Input: Query $q$, Graph $\mathcal{G}$, Policy $\pi_\theta$, Max Steps $T_{\max}$
Output: Candidate Answer Set $\mathcal{S}$
1: $\mathcal{S} \leftarrow []$; $\mathcal{H}_0 \leftarrow [q]$; $t \leftarrow 1$
2: while $t \le T_{\max}$ do
3:   Sample action $a_t \sim \pi_\theta(\cdot \mid \mathcal{H}_{t-1})$
4:   if $a_t = \text{GlobalSearch}(q_{\text{sub}}, k)$ then
5:     $o_t \leftarrow \text{BM25Search}(\mathcal{G}, q_{\text{sub}}, k)$
6:   elif $a_t = \text{NeighborhoodExploration}(v, q_{\text{sub}}, k)$ then
7:     $o_t \leftarrow \text{RankNeighbors}(\mathcal{G}, v, q_{\text{sub}}, k)$
8:   elif $a_t = \text{AddToAnswer}(v_1, \dots, v_m)$ then
9:     $\mathcal{S} \leftarrow \mathcal{S} \mathbin{\Vert} [v_1, \dots, v_m]$
10:    $o_t \leftarrow \text{"Nodes added"}$
11:  elif $a_t = \text{Finish}()$ then
12:    break
13:  end if
14:  $\mathcal{H}_t \leftarrow \mathcal{H}_{t-1} \mathbin{\Vert} (a_t, o_t)$
15:  $t \leftarrow t + 1$
16: end while
17: return $\mathcal{S}$
\end{lstlisting}
\end{figure}

\subsection{Multi-Agent Trajectories and Voting Rank Fusion}
Because graph exploration is inherently non-convex and subject to local search traps, ARK deploys $M$ independent agents in parallel. Each agent $m \in \{1, \dots, M\}$ generates an exploration trajectory $\tau_m = (a_1, o_1, a_2, o_2, \dots, a_{T_m})$ at temperature $T > 0$, producing an ordered candidate answer set $\mathcal{S}_m$.

To combine these heterogeneous answer sets, ARK applies \textbf{Voting Rank Fusion} (VRF). For any node $v \in \bigcup_{m=1}^M \mathcal{S}_m$, its fused score is computed via reciprocal rank weighting:
\begin{equation}
\label{eq:vrf}
\text{Score}_{\text{VRF}}(v) = \sum_{m=1}^M \mathbb{I}(v \in \mathcal{S}_m) \cdot \frac{1}{\text{rank}_m(v) + \kappa}
\end{equation}
where $\text{rank}_m(v) \in \{1, 2, \dots\}$ represents the 1-indexed position of node $v$ in agent $m$'s answer set $\mathcal{S}_m$, $\kappa \ge 0$ is a smoothing constant (default $\kappa = 60$), and $\mathbb{I}(\cdot)$ is the indicator function. The final candidate list $\mathcal{R}_K$ is obtained by sorting nodes in descending order of $\text{Score}_{\text{VRF}}(v)$, as detailed in Algorithm~\ref{alg:vrf}.

\begin{figure}[t]
\centering
\begin{lstlisting}[mathescape=true, caption={Voting Rank Fusion (VRF)}, label={alg:vrf}]
Input: Trajectory answer lists $\{\mathcal{S}_1, \dots, \mathcal{S}_M\}$, Top-$K$, Constant $\kappa$
Output: Ranked candidates $\mathcal{R}_K$
1: Initialize score map $\text{Scores} \leftarrow \{\}$
2: for $m = 1$ to $M$ do
3:   for $i = 1$ to $|\mathcal{S}_m|$ do
4:     $v \leftarrow \mathcal{S}_m[i]$
5:     $\text{Scores}[v] \leftarrow \text{Scores}.get(v, 0) + \frac{1}{i + \kappa}$
6:   end for
7: end for
8: Sort all $v \in \text{Scores}$ descending by $\text{Scores}[v]$
9: return Top-$K$ nodes from sorted list
\end{lstlisting}
\end{figure}

% =============================================================================
% 4. MODEL FINE-TUNING & DISTILLATION
% =============================================================================
\section{Trajectory Mining \& Model Distillation}
\label{sec:distill}

\begin{figure}[t]
\centering
\resizebox{\columnwidth}{!}{%
\begin{tikzpicture}[
    node distance=0.55cm and 0.5cm,
    proc/.style={rectangle, draw=navyblue, fill=blue!6, thick, rounded corners=3pt, align=center, inner sep=4pt, font=\scriptsize},
    databox/.style={rectangle, draw=darkred, fill=red!6, thick, rounded corners=3pt, align=center, inner sep=4pt, font=\scriptsize},
    modelbox/.style={rectangle, draw=darkgreen, fill=green!6, thick, rounded corners=3pt, align=center, inner sep=4pt, font=\scriptsize},
    arr/.style={-Latex, thick, navyblue}
]
\node[databox] (qa) {\textbf{STaRK-PrimeKG Questions}\\\textit{592 Train / 199 Val Queries}};

\node[proc, below left=0.55cm and -0.3cm of qa] (oracle) {\textbf{Oracle Shortest Path}\\\textit{Deterministic BFS}};
\node[proc, below right=0.55cm and -0.3cm of qa] (teacher) {\textbf{Gemini 3.6 Flash}\\\textit{5 Agents $\times$ Queries}};

\node[databox, below=0.45cm of oracle] (oracledata) {\textbf{Oracle Trajectories}\\592 Train / 199 Val\\(1.0 Ans/Traj)};
\node[databox, below=0.45cm of teacher] (rejection) {\textbf{Rejection Sampling}\\Keep $\text{Rec@20} > 0$ (Best/Q)\\123 Queries $\times 3 = 369$};

\node[modelbox, below=0.45cm of oracledata] (qwen4b) {\textbf{Qwen3-4B LoRA SFT}\\Full Precision bfloat16\\$\text{eval\_loss} = \mathbf{0.0752}$};
\node[modelbox, below=0.45cm of rejection] (qwen06b) {\textbf{Qwen3-0.6B QLoRA Distill}\\INT4 NF4 Quantized\\$\text{eval\_loss} = \mathbf{0.3680}$};

\draw[arr] (qa) -- (oracle);
\draw[arr] (qa) -- (teacher);
\draw[arr] (oracle) -- (oracledata);
\draw[arr] (teacher) -- (rejection);
\draw[arr] (oracledata) -- (qwen4b);
\draw[arr] (rejection) -- (qwen06b);

\end{tikzpicture}%
}
\caption{Dual Trajectory Mining and Distillation Flow. Left: Oracle shortest-path SFT for \texttt{Qwen3-4B}. Right: Gemini 3.6 Flash teacher exploration with Rejection Sampling for \texttt{Qwen3-0.6B}.}
\label{fig:distill_flow}
\end{figure}

To eliminate reliance on commercial APIs, we construct two complementary training regimes depicted in Figure~\ref{fig:distill_flow}.

\subsection{Oracle Shortest-Path Trajectories}
For supervised fine-tuning on ideal reasoning trajectories, we construct deterministic \textbf{Oracle} paths. Given question $q$ with source entity $v_{\text{src}}$ and target answer entity $v^*$, the Oracle computes the shortest relational path $v_{\text{src}} \xrightarrow{r_1} v_1 \xrightarrow{r_2} \dots \xrightarrow{r_k} v^*$ via Breadth-First Search (BFS) over $\mathcal{G}$.

The Oracle trajectory simulates an ideal agent executing:
\begin{enumerate}[leftmargin=*]
    \item $\text{GlobalSearch}(v_{\text{src}}, k)$ to discover the seed.
    \item Successive $\text{NeighborhoodExploration}$ calls along the shortest path.
    \item $\text{AddToAnswer}(v^*)$ followed by $\text{Finish}()$.
\end{enumerate}
This yields 592 training trajectories and 199 validation trajectories on PrimeKG, with an average target density of 1.0 answers per trajectory.

\subsection{Teacher Trajectory Mining \& Rejection Sampling}
While Oracle trajectories provide clean shortest paths, they lack exploratory behaviors (e.g., backtracking from dead ends or broad lexical searching when source entities are ambiguous). To capture robust exploration policies, we deploy Gemini~3.6~Flash with $M=5$ parallel agents across all 592 training and 199 validation queries, generating 2,840 training and 881 validation trajectories.

Because raw teacher trajectories contain exploratory failures, training directly on all traces degrades student precision. We apply \textbf{Rejection Sampling Fine-Tuning} (RFT):
\begin{enumerate}[leftmargin=*]
    \item \textit{Quality Filtering:} For each query $q$, we evaluate all $M=5$ candidate trajectories and retain only trajectories satisfying $\text{Recall@20}(\tau) > 0$.
    \item \textit{Optimal Selection:} Among surviving candidates, we select the single best trajectory $\tau^* = \arg\max_\tau \text{Recall@20}(\tau)$, using minimum trajectory step count $|\tau|$ as a tiebreaker.
    \item \textit{Data Balancing:} Each accepted trajectory is tripled ($3\times$) to maintain batch diversity.
\end{enumerate}
As detailed in Table~\ref{tab:dataset_stats}, rejection sampling retains 123 usable training queries (20.8\% acceptance) and 50 validation queries (25.1\% acceptance), yielding 369 training samples and 150 validation samples.

\begin{table}[t]
\centering
\small
\resizebox{\columnwidth}{!}{%
\begin{tabular}{lcccc}
\toprule
\textbf{Dataset Split} & \textbf{Total} & \textbf{Usable} & \textbf{Accept.} & \textbf{Samples} \\
\midrule
PrimeKG Train (Oracle) & 592 & 592 & 100\% & 592 \\
PrimeKG Val (Oracle) & 199 & 199 & 100\% & 199 \\
\midrule
Gemini Teacher Train & 592 & 123 & 20.8\% & 369 ($3\times$) \\
Gemini Teacher Val & 199 & 50 & 25.1\% & 150 ($3\times$) \\
\bottomrule
\end{tabular}%
}
\caption{Trajectory dataset statistics before and after Rejection Sampling filtering on STaRK-PrimeKG.}
\label{tab:dataset_stats}
\end{table}

\subsection{Assistant-Masked SFT Objective}
Let a trajectory sample be formatted as a token sequence $\mathbf{x} = (x_1, x_2, \dots, x_N)$, interleaved with system instructions, user queries, environment tool observations, and assistant thoughts/tool-calls.

We train student models using an \textbf{assistant-only masked cross-entropy loss}. We assign target labels $y_t = x_t$ if token $t$ is generated by the assistant, and $y_t = -100$ (ignored in gradient computation) otherwise:
\begin{equation}
\label{eq:sft_loss}
\resizebox{0.92\columnwidth}{!}{$
\displaystyle \mathcal{L}_{\text{SFT}}(\theta) = -\frac{1}{N_{\text{act}}} \sum_{t=1}^N \mathbb{I}(y_t \ne -100) \log P_\theta(x_t \mid x_{<t})
$}
\end{equation}
where $N_{\text{act}} = \sum_{t=1}^N \mathbb{I}(y_t \ne -100)$ denotes the total number of active assistant tokens. This ensures the model learns only the policy distribution $P(a_t \mid h_t)$ without expending model capacity on predicting graph environment responses.

\subsection{Parameter-Efficient Adaptation Architectures}
We configure two distinct model scale regimes as summarized in Table~\ref{tab:full_hypers}:

\begin{table}[t]
\centering
\small
\resizebox{\columnwidth}{!}{%
\begin{tabular}{lcc}
\toprule
\textbf{Hyperparameter} & \textbf{Qwen3-4B FT} & \textbf{Qwen3-0.6B Dist.} \\
\midrule
Base Model & \texttt{Qwen/Qwen3-4B} & \texttt{Qwen/Qwen3-0.6B} \\
Precision & bfloat16 & INT4 (NF4) \\
LoRA Rank ($r$) & 32 & 32 \\
LoRA Alpha ($\alpha$) & 64 & 64 \\
LoRA Dropout & 0.1 & 0.1 \\
Target Modules & All linear & All linear \\
Epochs & 3 & 5 \\
Effective Batch Size & 16 & 16 \\
Learning Rate & $1 \times 10^{-4}$ & $5 \times 10^{-5}$ \\
Warmup Steps & 50 & 30 \\
Max Sequence Length & 4,096 & 4,096 \\
Optimizer & AdamW & AdamW \\
Loss Formulation & Assistant-only & Assistant-only \\
Rejection Sampling & No (Oracle) & Yes ($\text{Rec@20}>0$) \\
\bottomrule
\end{tabular}%
}
\caption{Full hyperparameter configuration for student models.}
\label{tab:full_hypers}
\end{table}

\paragraph{1. \texttt{Qwen3-4B} Oracle Fine-Tuning:}
Trained on Oracle paths using full-precision bfloat16 Low-Rank Adaptation (LoRA) \citep{hu2022lora} with rank $r=32$, $\alpha=64$, and dropout $0.1$. LoRA adapters are injected into all attention projections (\texttt{q\_proj, k\_proj, v\_proj, o\_proj}) and feed-forward projections (\texttt{gate\_proj, up\_proj, down\_proj}). Optimization is performed over 3 epochs (111 steps) with effective batch size 16 (batch size 4, gradient accumulation 4) and learning rate $1\times 10^{-4}$ under AdamW.

\paragraph{2. \texttt{Qwen3-0.6B} Gemini Distillation:}
To enable deployment on edge GPUs with $<2\,\text{GB}$ VRAM, \texttt{Qwen3-0.6B} is quantized to 4-bit NormalFloat (NF4) with double quantization via QLoRA \citep{dettmers2023qlora}. Adapters ($r=32, \alpha=64$) are attached to all linear layers. The model is trained over 5 epochs (120 steps) with effective batch size 16 (batch size 1, gradient accumulation 16) and learning rate $5\times 10^{-5}$.

% =============================================================================
% 5. EXPERIMENTS AND RESULTS
% =============================================================================
\section{Experiments and Empirical Results}
\label{sec:experiments}

\subsection{Experimental Setup and Benchmark}
We evaluate our methods on the \textbf{STaRK-PrimeKG} benchmark \citep{wu2024stark,chandak2023primekg}. PrimeKG integrates 20 biomedical databases into a heterogeneous graph comprising $|\mathcal{V}| = 129,375$ nodes across 10 biological types (genes, diseases, drugs, phenotypes, pathways) and $|\mathcal{E}| = 8,110,634$ directed edges across 29 relationship types, with a high average node degree of 125.2.

\paragraph{Evaluation Metrics.}
Given ground-truth target entities $\mathcal{V}^*$ and retrieved top-$K$ candidates $\mathcal{R}_K$, we report standard retrieval metrics:
\begin{itemize}[leftmargin=*]
    \item \textbf{Hit@$K$ ($K \in \{1, 5, 10\}$):} $\text{Hit@}K = \mathbb{I}(|\mathcal{V}^* \cap \mathcal{R}_K| > 0)$, checking whether at least one target is retrieved.
    \item \textbf{Recall@$K$ ($K \in \{10, 20\}$):} $\text{Rec@}K = \frac{|\mathcal{V}^* \cap \mathcal{R}_K|}{|\mathcal{V}^*|}$, computing the proportion of targets recovered.
    \item \textbf{MRR (Mean Reciprocal Rank):} $\text{MRR} = \frac{1}{\min \{ \text{rank}(v) \mid v \in \mathcal{V}^* \cap \mathcal{R}_K \}}$, measuring ranking quality.
\end{itemize}

\begin{table*}[t]
\centering
\small
\resizebox{\textwidth}{!}{%
\begin{tabular}{llcccccccc}
\toprule
\textbf{Model / Method} & \textbf{Split} & \textbf{Agents} & $n$ & \textbf{Hit@1} & \textbf{Hit@5} & \textbf{Hit@10} & \textbf{Rec@10} & \textbf{Rec@20} & \textbf{MRR} \\
\midrule
\textit{Oracle (Upper Bound)} & Val & 1 & 199 & 0.5879 & 0.5879 & 0.5879 & 0.4511 & 0.4511 & 0.5879 \\
\textit{Oracle (Upper Bound)} & Train & 1 & 592 & 0.5507 & 0.5507 & 0.5507 & 0.4091 & 0.4091 & 0.5507 \\
\midrule
\textbf{Gemini 3.6 Flash (Teacher)} & Val & 5 & 199 & 0.2060 & 0.2312 & 0.2312 & 0.1788 & 0.1904 & 0.2180 \\
\textbf{Gemini 3.6 Flash (Teacher)} & Train & 5 & 592 & 0.1537 & 0.1909 & 0.1943 & 0.1473 & 0.1574 & 0.1693 \\
\midrule
\textbf{Qwen3-4B LoRA SFT (Ours)} & Val & 1 & 20 & \textbf{0.2000} & \textbf{0.2000} & \textbf{0.2000} & \textbf{0.2000} & \textbf{0.2000} & \textbf{0.2000} \\
\textbf{Qwen3-0.6B QLoRA Distilled (Ours)} & Val & 3 & 199 & \multicolumn{6}{c}{\textit{Evaluation in progress ($\text{eval\_loss} = 0.3680$)}} \\
\bottomrule
\end{tabular}%
}
\caption{Retrieval performance on STaRK-PrimeKG across Oracle, Teacher, and Fine-Tuned student models. Multi-agent ensembles use Voting Rank Fusion ($\kappa=60$).}
\label{tab:main_results}
\end{table*}

\subsection{Main Retrieval Results}
Table~\ref{tab:main_results} reports retrieval performance across models.

\paragraph{1. The Oracle Ceiling:}
The Oracle shortest-path achieves $\text{Hit@1} = 0.5879$ and $\text{MRR} = 0.5879$ on the validation set. This represents the empirical upper bound achievable when an agent follows the exact shortest relational path connecting query entities to answers. The gap between Oracle and LLM agents illustrates the inherent difficulty of navigating dense graphs without a global topological map.

\paragraph{2. Teacher Baseline Performance:}
Gemini 3.6 Flash with $M=5$ parallel agents and Voting Rank Fusion achieves $\text{Hit@1} = 0.2060$ and $\text{MRR} = 0.2180$ on the validation set. Ensembling 5 stochastic trajectories boosts $\text{Hit@5}$ from $0.2060$ to $0.2312$ (+2.52 percentage points), confirming the benefit of multi-agent rank aggregation.

\paragraph{3. Student Parity with Single Agent:}
Remarkably, our fine-tuned \texttt{Qwen3-4B} model operating with only a \textbf{single agent} ($M=1$) achieves $\text{Hit@1} = 0.2000$ on the validation sample ($n=20$), matching the multi-agent ensemble of the teacher ($0.2060$). This demonstrates that training on structured shortest-path trajectories instills efficient navigation heuristics, allowing a 4B parameter model to achieve competitive retrieval accuracy without multi-agent overhead.

\begin{table}[t]
\centering
\small
\resizebox{\columnwidth}{!}{%
\begin{tabular}{lcccc}
\toprule
\textbf{Step} & \textbf{Epoch} & \textbf{Train Loss} & \textbf{Eval Loss} & \textbf{Learning Rate} \\
\midrule
\multicolumn{5}{l}{\textit{\textbf{Qwen3-4B Oracle SFT}} (Total: 111 steps, 3 epochs)} \\
10 & 0.27 & 1.9128 & -- & $1.8 \times 10^{-5}$ \\
30 & 0.81 & 0.2528 & -- & $5.8 \times 10^{-5}$ \\
50 & 1.35 & 0.0909 & -- & $9.8 \times 10^{-5}$ \\
80 & 2.16 & 0.0693 & -- & $5.2 \times 10^{-5}$ \\
110 & 2.97 & 0.0651 & -- & $3.0 \times 10^{-6}$ \\
\textbf{111} & \textbf{3.00} & \textbf{0.0642} & $\mathbf{0.0752}$ & -- \\
\midrule
\multicolumn{5}{l}{\textit{\textbf{Qwen3-0.6B Gemini Distillation}} (Total: 120 steps, 5 epochs)} \\
20 & 0.83 & 0.6472 & -- & $3.2 \times 10^{-5}$ \\
50 & 2.08 & 0.3226 & 0.3916 & $3.9 \times 10^{-5}$ \\
80 & 3.33 & 0.3132 & -- & $2.3 \times 10^{-5}$ \\
100 & 4.17 & 0.2741 & 0.3698 & $1.2 \times 10^{-5}$ \\
\textbf{120} & \textbf{5.00} & \textbf{0.2644} & $\mathbf{0.3680}$ & $1.0 \times 10^{-6}$ \\
\bottomrule
\end{tabular}%
}
\caption{Training and validation loss progression during fine-tuning of \texttt{Qwen3-4B} and distillation of \texttt{Qwen3-0.6B}.}
\label{tab:training_loss}
\end{table}

\subsection{Training Dynamics and Convergence Analysis}
Table~\ref{tab:training_loss} presents the training progression for both student models.

\paragraph{Rapid Convergence of \texttt{Qwen3-4B}:}
For \texttt{Qwen3-4B}, the training loss drops precipitously from $1.9128$ at step 10 to $0.0909$ at step 50 during the warmup phase, stabilizing at $\sim 0.065$ in subsequent epochs. The final validation loss reaches $\mathbf{0.0752}$ at step 111. Because Oracle trajectories follow strict shortest-path syntax, the model rapidly internalizes the state-action transition distribution.

\paragraph{Generalization in \texttt{Qwen3-0.6B} Distillation:}
For \texttt{Qwen3-0.6B}, training loss decreases steadily from $2.3703$ to $0.2644$ across 5 epochs. Validation loss improves monotonically ($0.3916 \to 0.3698 \to \mathbf{0.3680}$), confirming that rejection sampling effectively eliminates noisy trajectories and prevents student overfitting.

\subsection{Tool Ablation: Lexical vs. Dense Global Search}
A core architectural premise of ARK is employing lexical BM25 rather than dense embeddings for $\text{GlobalSearch}$. To validate this design choice, we compare BM25 against dense cosine retrieval using \texttt{text-embedding-3-large} over 1,640 test queries from the STaRK-Amazon benchmark.

As reported in Table~\ref{tab:tool_ablation}, lexical BM25 search substantially outperforms dense embedding retrieval across all metrics under identical agent configurations ($\text{Hit@1} = 0.5640$ vs. $0.4713$, +9.27 percentage points; $\text{MRR} = 0.6544$ vs. $0.5745$). In heterogeneous graphs, query entity names require exact lexical matching rather than semantic paraphrase mapping, where dense embeddings frequently retrieve spurious semantic neighbors.

\begin{table}[t]
\centering
\small
\resizebox{\columnwidth}{!}{%
\begin{tabular}{lcccc}
\toprule
\textbf{Search Mode} & \textbf{Hit@1} & \textbf{Hit@5} & \textbf{Rec@20} & \textbf{MRR} \\
\midrule
Dense Embeddings & 0.4713 & 0.6994 & 0.5813 & 0.5745 \\
\textbf{Lexical BM25 (ARK)} & \textbf{0.5640} & \textbf{0.7622} & \textbf{0.6058} & \textbf{0.6544} \\
\midrule
\textit{$\Delta$ Gain (BM25)} & \textbf{+9.27\%} & \textbf{+6.28\%} & \textbf{+2.45\%} & \textbf{+7.99\%} \\
\bottomrule
\end{tabular}%
}
\caption{Ablation of Global Search retrieval modality using \texttt{GPT-4.1} explorer on STaRK-Amazon test split ($n=1,640$).}
\label{tab:tool_ablation}
\end{table}

\subsection{Computational Footprint and Serving Efficiency}
Table~\ref{tab:compute_efficiency} compares the operational footprint of frontier API models against our fine-tuned local models.

\begin{table}[t]
\centering
\small
\resizebox{\columnwidth}{!}{%
\begin{tabular}{lccc}
\toprule
\textbf{Model} & \textbf{Params} & \textbf{VRAM (Est.)} & \textbf{Latency / Query} \\
\midrule
Gemini 3.6 Flash & Frontier & API Cloud & $\sim 4.8\,\text{s}$ (5 agents) \\
\texttt{Qwen3-4B} FT & 4B & $\sim 3.5\,\text{GB}$ (INT4) & $\sim 0.9\,\text{s}$ (1 agent) \\
\texttt{Qwen3-0.6B} Dist. & 0.6B & $\sim 1.2\,\text{GB}$ (INT4) & $\sim 0.3\,\text{s}$ (1 agent) \\
\bottomrule
\end{tabular}%
}
\caption{Serving resource footprint and latency comparison.}
\label{tab:compute_efficiency}
\end{table}

While proprietary APIs require network round-trips and incur recurring token fees, our distilled \texttt{Qwen3-0.6B} model occupies only $\sim 1.2\,\text{GB}$ of VRAM under INT4 quantization, executing single-agent queries in under $300\,\text{ms}$ on commodity hardware.

% =============================================================================
% 6. QUALITATIVE CASE STUDY
% =============================================================================
\section{Qualitative Analysis and Case Study}
\label{sec:casestudy}

To illustrate the trajectory execution flow, Figure~\ref{fig:trace_example} shows an actual trace from the validation set of STaRK-PrimeKG.

\begin{figure}[h]
\centering
\begin{lstlisting}[language=Python]
# Query: "Which molecular target is modulated by Drug X
# and associated with cardiomyopathy Y?"

# Step 1: Global Discovery
Action: global_search(subquery="Drug X", k=3)
Observation: [Node 1042: "Drug X (Small Molecule)"]

# Step 2: Neighborhood Relational Expansion
Action: neighborhood_exploration(node_id=1042,
                                 subquery="cardiomyopathy", k=5)
Observation: [
  (Node 8521, "Target Gene ABC", "drug_protein"),
  (Node 9140, "Enzyme DEF", "drug_protein")
]

# Step 3: Candidate Verification & Add
Action: add_to_answer(node_ids=[8521])
Observation: "Node 8521 added to candidate set."

# Step 4: Trajectory Termination
Action: finish()
Observation: "Exploration terminated."
\end{lstlisting}
\caption{Representative ARK trajectory trace demonstrating initial lexical seed discovery followed by relational neighborhood expansion.}
\label{fig:trace_example}
\end{figure}

The agent first initiates $\text{global\_search}$ to locate the specific drug entity node \texttt{1042}. Next, instead of blindly expanding all adjacent relationships, it invokes $\text{neighborhood\_exploration}$ filtered by the disease keyword \textit{``cardiomyopathy''}, identifying the candidate protein \texttt{8521}. Finally, it adds the node to its answer buffer and concludes via $\text{finish}()$.

% =============================================================================
% 7. CONCLUSION
% =============================================================================
\section{Conclusion}
\label{sec:conclusion}
In this paper, we investigated the ARK framework for autonomous multi-agent knowledge graph exploration and addressed the challenge of small language model specialization. By developing a dual trajectory mining pipeline comprising Oracle shortest paths and Rejection-Sampled teacher trajectories, we trained \texttt{Qwen3-4B} and \texttt{Qwen3-0.6B} to navigate heterogeneous biomedical graphs. Empirical evaluation on STaRK-PrimeKG proves that lightweight fine-tuned models can match 5-agent frontier teacher ensembles with a single agent, achieving high sample efficiency and fast convergence while drastically reducing inference latency and cost.

% =============================================================================
% REFERENCES
% =============================================================================
\bibliography{references}

\end{document}
